\documentclass[letterpaper]{article}
\usepackage[submission]{aaai2027}
\usepackage[hyphens]{url}
\usepackage{graphicx}
\urlstyle{rm}
\def\UrlFont{\rm}
\usepackage{natbib}
\usepackage{caption}
\frenchspacing
\usepackage{booktabs}
\pdfinfo{
/TemplateVersion (2027.1)
}
\setcounter{secnumdepth}{0}

\title{PaperToSkill: Turning Research Papers into Portable Agent Skills}
\author{
    Anonymous Authors
}
\affiliations{
    Anonymous Institution\\
    anonymous@example.com
}

\begin{document}

\maketitle

\begin{abstract}
Many LLM and agent papers introduce workflows that could help ordinary users and
agent builders, but those workflows are usually described for human reading
rather than direct reuse. Summaries compress the idea, but often discard the
procedural details, validation checks, assumptions, and failure branches needed
to apply the method. We introduce PaperToSkill, a paper-to-skill conversion
workflow that turns source-anchored paper notes into compact, human-editable
agent skills. Each skill records the paper's contribution, use conditions,
workflow, validation checks, failure cases, transfer notes, and source anchors.
In a four-paper benchmark over AI Scientist-v2, Reflexion, AIDE, and
Toolformer, PaperToSkill generates skills that score 20/20 on deterministic
structural rubrics, preserve more deterministic operational coverage than
generic-summary and abstract-only baselines, remain under a 1200-word
compactness budget, and substantially reduce deterministic input-token proxy
size relative to full extracted papers. These results support PaperToSkill as a
reproducible conversion layer from curated paper notes to portable skills. All
four live-transfer saved-response sets are collected and scored under a
deterministic output-contract evaluator; these saved response files are not
human semantic fidelity or downstream execution-outcome evidence. A first
single-run real-reuse stress test over eight AIDE, SWE-agent, Reflexion, and
SnapATAC2 paper-task rows is complete, but the outcome is mixed rather than
confirmatory: AIDE-T2 and one SWE task favor PaperToSkill, AIDE-T1 and
Reflexion are solved by both conditions, and SWE-T1 plus SnapATAC2 expose
patch-application and artifact-completion failures. We therefore treat these
rows as downstream boundary evidence, not aggregate effectiveness.
\end{abstract}

\section{Introduction}

Using LLMs is becoming a practical skill for researchers, builders, and
non-expert users. At the same time, many of the best ways to use LLM agents are
introduced in papers whose methods are difficult to operationalize. A paper may
describe a search policy, memory loop, validation stage, or interface contract,
but a user who wants to reuse the contribution must translate prose into an
agent workflow by hand.

PaperToSkill studies a simple hypothesis: papers can be converted into compact,
editable skills that preserve enough procedural knowledge for agents and users
to reuse the main contribution. We use ``skill'' to mean a natural-language
operational artifact with an entry-point \texttt{SKILL.md}, similar to formats
used by agent harnesses that load task-specific instructions. Unlike a summary,
a skill is not only explanatory. It should specify when to use the method,
which inputs are required, which steps to execute, how to validate progress,
which failures to watch for, and how to adapt the instructions across
harnesses.

This paper takes an artifact-first view. PaperToSkill is not presented as a
fully automatic PDF understanding system. The current implementation converts
curated, source-anchored paper notes into skills and source maps. This narrow
scope is intentional: it lets us evaluate whether the conversion layer
preserves operational knowledge before claiming end-to-end PDF automation.

Our current contributions are: (1) a paper-to-skill schema for operationalizing
research contributions; (2) a deterministic extraction scaffold that emits
\texttt{SKILL.md} plus source maps; (3) deterministic extracted-text-to-note
scaffolds evaluated on Toolformer and AIDE; (4) a four-paper benchmark using AI
Scientist-v2, Reflexion, AIDE, and Toolformer; (5) deterministic/offline
evaluations for structure, context coverage, compactness, source grounding, and
transfer readiness; and (6) a first single-run real-reuse stress test that
records both positive and failed Summary-vs-PaperToSkill outcomes without
promoting them into an aggregate downstream-success claim.

\section{Related Work}

PaperToSkill is motivated by automated research and agent workflow systems.
AI Scientist-v2 uses agentic tree search and staged experiment management for
automated scientific discovery~\cite{lu2025aiscientistv2}. AIDE frames ML
engineering as search over code solutions~\cite{aide2025}. Reflexion stores
verbal feedback in episodic memory to improve later attempts without weight
updates~\cite{shinn2023reflexion}. Toolformer shows how language models can
generate and filter tool-use examples for API calls~\cite{schick2023toolformer}.
Skill-library and interface work such as Voyager and SWE-agent also demonstrate
that reusable skills and agent-facing interfaces can determine whether a method
transfers beyond its original implementation~\cite{wang2023voyager,yang2024sweagent}.
Paper2Agent is the closest concurrent system: it converts papers and associated
codebases into MCP servers with tools, resources, prompts, tests, and an
interactive paper-agent layer~\cite{miao2025paper2agent}. AgenticSciML shows a
related use of specialized agents, structured debate, method memory, and
evolutionary search for scientific machine-learning discovery~\cite{jiang2026agenticsciml}.

These systems show that papers can become active computational artifacts.
PaperToSkill targets a different conversion layer: compact, human-editable
skills that preserve operational knowledge, validation checks, failure branches,
transfer notes, and source boundaries without requiring an MCP server or a
runnable paper codebase.

\section{Method}

PaperToSkill represents each converted paper as a skill package. The main file
is \texttt{SKILL.md}; optional references include a source map that links
generated skill bullets to source note sections and line spans. The schema
contains identity, central contribution, use conditions, required artifacts,
workflow steps, validation checks, failure cases, transfer notes, and source
notes that distinguish source-backed material from inferred guidance.

The deterministic extractor reads a source-anchored paper note, normalizes
Markdown sections and list items, collects candidates from abstract, methods,
experiments, limitations, and transfer sections, then emits a compact skill.
Candidate limits keep the artifact within a practical context budget while
preserving method, validation, and failure coverage. The extractor also writes a
\texttt{references/source\_map.json} file so later audits can inspect which
source section supports each generated instruction.

The text-to-note scaffold is an earlier preprocessing step. It operates on
newline-numbered extracted text, detects broad paper regions, selects line
windows for methods, experiments, and limitations, and preserves line anchors.
For two-column extracted text, it keeps raw line spacing and selects the
keyword-bearing column while retaining the original source line range. The
output is explicitly marked as an audit scaffold.

\section{Experimental Setup}

We evaluate four real agent-method papers: AI Scientist-v2, Reflexion, AIDE,
and Toolformer. For each paper, we create a curated source-anchored note from
the extracted paper text, generate a PaperToSkill skill, and compare it against
two context baselines: a generic summary and an abstract-only context. For
transfer ablation, we compare the full skill against the same skill with
\texttt{Transfer Notes} removed and against the generic summary.

The metrics are deterministic: skill rubric score, context coverage,
source-span validation, offline transfer-readiness, word count under a
1200-word compactness budget, character-based token/cost proxy, and a local
tokenizer-aware input proxy using \texttt{o200k\_base}. For saved
model-ablation responses, we additionally compute a local output-token proxy.
These metrics are reproducible gates. They do not replace live agent execution,
real billing records, invoice evidence, success-per-dollar accounting, or
completed human fidelity annotation. The current package uses local token
accounting over input-token and saved-response output-token proxies. We also
prepare model-ablation prompt packets for Claude Opus 4.8, a GPT-family slot
requested as GPT 5.5, and a DeepSeek slot, and include a protocol-aware runner
and response evaluator. The current two-case protocol has 6/6 saved and scored
rows across Claude, GPT-family, and DeepSeek slots. This is saved-response
output-contract evidence only; it is not live downstream task success,
provider billing, or a broad model-quality result. Provider/model availability
failures are logged separately from model-quality results.

The primary downstream experiment is the real-reuse protocol in
Table~\ref{tab:real-reuse-main}. It contains eight original-style
\emph{paper-task} rows across AIDE, SWE-agent, Reflexion, and SnapATAC2. Each
row fixes the source paper, task input/output shape, metric family, reported
reference boundary, and the two main conditions: Summary and PaperToSkill. At
this revision, all eight rows have one GPT-family run with scored Summary and
PaperToSkill outputs. The AIDE rows use an official Kaggle Spaceship Titanic
training file with a locked local validation split; the SnapATAC2 rows use
prepared official miniature fixtures and are scored local dry-run tasks rather
than full paper reproductions. These rows complete the first single-run pass
over the real-reuse table, but they do not by themselves establish aggregate
advantage over Summary.

We also instantiate a small Full Excerpt sanity scaffold over the
pre-registered AIDE-T1, SWE-T1, and SNAP-T1 subset
(Table~\ref{tab:full-excerpt-sanity}). This table is separate from the main
real-reuse comparison: it records the existing Summary/PaperToSkill scores and
local context-token proxies, and the matched Full Excerpt runs remain
auxiliary sanity evidence rather than a main baseline.

\begin{table*}[t]
\centering
\caption{Main real-reuse experiment table. Filled scores come from local raw rows; future reruns or additional rows may be added. Original-paper scores are reported references unless the same data, input/output, metric, budget, and runtime setting are reproduced locally.}
\label{tab:real-reuse-main}
\resizebox{\textwidth}{!}{%
\begin{tabular}{llllllllll}
\toprule
Task & Paper & Domain & Input & Output & Metric & Reference & Summary & PaperToSkill & Status \\
\midrule
AIDE-T1 & AIDE & ML engineering & Kaggle-style dataset + metric & Runnable solution/submission & validation\_score & Reported AIDE ref. & 0.816 & 0.817 & Scored (GPT-family) \\
AIDE-T2 & AIDE & ML engineering & Weak ML script + feedback & Improved script + trajectory & best\_node\_score & Reported AIDE ref. & 0.000 & 0.826 & Scored (GPT-family) \\
SWE-T1 & SWE-agent & Software engineering & Repo issue + tests & Patch + test log & resolved & Reported SWE-agent ref. & 0.000 & 0.000 & Scored (GPT-family) \\
SWE-T2 & SWE-agent & Software engineering & Failing test + repo & Focused patch + verification & tests\_passed & Reported SWE-agent ref. & 0.000 & 1.000 & Scored (GPT-family) \\
REF-T1 & Reflexion & Reasoning / QA & Multi-hop QA + feedback & Final answer + reflection trace & exact\_match\_or\_f1 & Reported Reflexion ref. & 1.000 & 1.000 & Scored (GPT-family) \\
REF-T2 & Reflexion & Decision / programming & Failed attempt + checker feedback & Corrected second attempt & second\_attempt\_success & Reported Reflexion ref. & 1.000 & 1.000 & Scored (GPT-family) \\
SNAP-T1 & SnapATAC2 & Single-cell omics & Small single-cell dataset & Pipeline + embedding artifacts & runtime\_memory\_quality & Reported SnapATAC2 ref. & 0.000 & 0.500 & Scored (GPT-family) \\
SNAP-T2 & SnapATAC2 & Single-cell omics & Single-cell labels/proxy task & Clustering/marker artifacts & ari\_nmi\_runtime\_memory & Reported SnapATAC2 ref. & 0.200 & 0.400 & Scored (GPT-family) \\
\bottomrule
\end{tabular}%
}
\end{table*}

\begin{table*}[t]
\centering
\caption{Real-reuse failure-boundary analysis derived from the same scored raw rows as Table~\ref{tab:real-reuse-main}. This table explains first-pass boundary modes; it does not add new task-success evidence.}
\label{tab:real-reuse-failure-analysis}
\resizebox{\textwidth}{!}{%
\begin{tabular}{lllll}
\toprule
Task & Summary Outcome & PaperToSkill Outcome & Boundary Mode & Contract Implication \\
\midrule
AIDE-T1 & 0.816; success & 0.817; success & Solved by both & harder task slice or stricter scorer \\
AIDE-T2 & 0.000; timeout & 0.826; success & PaperToSkill-only success & preserve patch/tool-use constraints \\
SWE-T1 & 0.000; patch apply failed & 0.000; patch apply failed & Patch application & patch-format and apply-check contract \\
SWE-T2 & 0.000; patch apply failed & 1.000; success & PaperToSkill-only success & preserve patch/tool-use constraints \\
REF-T1 & 1.000; success & 1.000; success & Solved by both & harder task slice or stricter scorer \\
REF-T2 & 1.000; success & 1.000; success & Solved by both & harder task slice or stricter scorer \\
SNAP-T1 & 0.000; artifact parse failed & 0.500; artifact incomplete & Artifact completion & required artifact and metric manifest \\
SNAP-T2 & 0.200; artifact incomplete & 0.400; artifact incomplete & Artifact completion & required artifact and metric manifest \\
\bottomrule
\end{tabular}%
}
\end{table*}

\begin{table*}[t]
\centering
\caption{SWE-T1 shared-source-context follow-up. The first-pass SWE-T1 row remains the main-table evidence; this paired follow-up exposes the same locked SQLFluff source slice to both conditions and is diagnostic rather than a replacement.}
\label{tab:swe-t1-source-context-followup}
\resizebox{\textwidth}{!}{%
\begin{tabular}{llllllllllll}
\toprule
Task & Condition & First Run & First Score & First Failure & First Patch & Follow-up Run & Follow-up Score & Follow-up Patch & Follow-up Test & Follow-up Failure & Interpretation \\
\midrule
SWE-T1 & Summary & phase97\_gpt\_swe\_t1\_summary\_retry & 0.000 & patch\_apply\_failed & No & phase107\_gpt\_swe\_t1\_source\_context\_followup & 0.000 & Yes & No & test\_command\_failed & Source context fixed patch application; hidden test still failed \\
SWE-T1 & PaperToSkill & phase97\_gpt\_swe\_t1\_real\_reuse & 0.000 & patch\_apply\_failed & No & phase107\_gpt\_swe\_t1\_source\_context\_followup & 0.000 & Yes & No & test\_command\_failed & Source context fixed patch application; hidden test still failed \\
\bottomrule
\end{tabular}%
}
\end{table*}

\begin{table*}[t]
\centering
\caption{SWE-T1 issue-aligned follow-up. The first-pass SWE-T1 row remains the main-table evidence; this diagnostic uses the pre-registered issue-aligned scorer and shows that both Summary and PaperToSkill pass under the revised contract, so it is not evidence of PaperToSkill advantage.}
\label{tab:swe-t1-issue-aligned-followup}
\resizebox{\textwidth}{!}{%
\begin{tabular}{llllllllllll}
\toprule
Task & Condition & Main Run & Main Score & Phase107 Run & Phase107 Score & Phase107 Test & Issue Run & Issue Score & Issue Test & Issue Test Patch & Interpretation \\
\midrule
SWE-T1 & Summary & phase97\_gpt\_swe\_t1\_summary\_retry & 0.000 & phase107\_gpt\_swe\_t1\_source\_context\_followup & 0.000 & No & phase110\_gpt\_swe\_t1\_issue\_aligned\_followup & 1.000 & Yes & No & Issue-aligned scorer passes both conditions; no PaperToSkill advantage \\
SWE-T1 & PaperToSkill & phase97\_gpt\_swe\_t1\_real\_reuse & 0.000 & phase107\_gpt\_swe\_t1\_source\_context\_followup & 0.000 & No & phase110\_gpt\_swe\_t1\_issue\_aligned\_followup & 1.000 & Yes & No & Issue-aligned scorer passes both conditions; no PaperToSkill advantage \\
\bottomrule
\end{tabular}%
}
\end{table*}

\begin{table*}[t]
\centering
\caption{SnapATAC2 executable-artifact follow-up. The main SNAP rows in Table~\ref{tab:real-reuse-main} remain unchanged; this paired diagnostic executes a pre-registered controlled scaffold over the same miniature fixtures to test whether concrete artifacts plus runtime/memory records satisfy the existing scorer.}
\label{tab:snapatac2-executable-followup}
\resizebox{\textwidth}{!}{%
\begin{tabular}{lllllll}
\toprule
Task & Condition & Run & Score & Success & Runtime Seconds & Peak Memory MB \\
\midrule
SNAP-T1 & summary & phase108\_snapatac2\_executable\_artifact\_followup & 1.0 & True & 1.296368 & 26.378289 \\
SNAP-T1 & papertoskill & phase108\_snapatac2\_executable\_artifact\_followup & 1.0 & True & 1.209319 & 26.338634 \\
SNAP-T2 & summary & phase108\_snapatac2\_executable\_artifact\_followup & 1.0 & True & 1.75799 & 24.648565 \\
SNAP-T2 & papertoskill & phase108\_snapatac2\_executable\_artifact\_followup & 1.0 & True & 1.756513 & 24.758577 \\
\bottomrule
\end{tabular}%
}
\end{table*}

\begin{table*}[t]
\centering
\caption{SnapATAC2 SNAP-T1 executable-candidate follow-up. The main SNAP rows in Table~\ref{tab:real-reuse-main} remain unchanged; this diagnostic executes model-generated phase112 candidate scripts under the pre-registered executable-candidate contract and shows contract closure for both conditions, not PaperToSkill advantage.}
\label{tab:snapatac2-executable-candidate-followup}
\resizebox{\textwidth}{!}{%
\begin{tabular}{lllllll}
\toprule
Task & Condition & Run & Score & Success & Runtime Seconds & Peak Memory MB \\
\midrule
SNAP-T1 & summary & phase112\_gpt\_snapatac2\_executable\_candidate\_t1 & 1.0 & True & 9.362002 & 0.0 \\
SNAP-T1 & papertoskill & phase112\_gpt\_snapatac2\_executable\_candidate\_t1 & 1.0 & True & 2.010309 & 164.615443 \\
\bottomrule
\end{tabular}%
}
\end{table*}

\begin{table*}[t]
\centering
\caption{Auxiliary Full Excerpt sanity check over the pre-registered subset. Scores come from matched real-reuse rows where available; token columns are local whitespace context proxies, not provider billing or output-token costs.}
\label{tab:full-excerpt-sanity}
\resizebox{\textwidth}{!}{%
\begin{tabular}{llllllllll}
\toprule
Task & Paper & Metric & Summary & PaperToSkill & Full Excerpt & Summary Tokens & PaperToSkill Tokens & Full Excerpt Tokens & Status \\
\midrule
AIDE-T1 & AIDE & validation\_score & 0.816 & 0.817 & 0.000 & 121 & 878 & 7366 & Scored (GPT-family) \\
SWE-T1 & SWE-agent & resolved & 0.000 & 0.000 & 0.000 & 89 & 1173 & 42048 & Scored (GPT-family) \\
SNAP-T1 & SnapATAC2 & runtime\_memory\_quality & 0.000 & 0.500 & 0.250 & 72 & 1069 & 10297 & Scored (GPT-family) \\
\bottomrule
\end{tabular}%
}
\end{table*}

\begin{table*}[t]
\centering
\caption{Main deterministic/offline PaperToSkill results. Coverage scores are task-rubric scores; transfer readiness is an offline artifact-readiness metric.}
\label{tab:main-results}
\begin{tabular}{lrrrrrrrr}
\toprule
Paper & Rubric & Skill & Summary & Abstract & $\Delta$Sum. & Transfer & Support & Words \\
\midrule
AI Scientist-v2 & 20/20 & 7.867/9 & 1.733/9 & 1.2/9 & 6.134 & 10/10 & 0.938 & 782 \\
Reflexion & 20/20 & 8.267/9 & 3.483/9 & 2.533/9 & 4.784 & 10/10 & 1.000 & 479 \\
AIDE & 20/20 & 9.1/10 & 1.916/10 & 1.333/10 & 7.184 & 10/10 & 1.000 & 927 \\
Toolformer & 20/20 & 8.9/10 & 2.5/10 & 1.534/10 & 6.400 & 10/10 & 1.000 & 943 \\
\bottomrule
\end{tabular}
\end{table*}

\begin{table}[t]
\centering
\caption{Transfer-note ablation. Removing only \texttt{Transfer Notes} lowers offline readiness across all curated skills.}
\label{tab:transfer-ablation}
\begin{tabular}{lrr}
\toprule
Paper & Full & No Transfer \\
\midrule
AI Scientist-v2 & 10/10 & 7.6/10 \\
Reflexion & 10/10 & 7.6/10 \\
AIDE & 10/10 & 7.6/10 \\
Toolformer & 10/10 & 7.6/10 \\
\bottomrule
\end{tabular}
\end{table}

\begin{table*}[t]
\centering
\caption{Local tokenizer-aware context-size proxy using the \texttt{o200k\_base} tokenizer. Counts are not provider bills, output-token costs, or success-per-dollar measurements.}
\label{tab:cost-proxy}
\begin{tabular}{lrrrr}
\toprule
Paper & Full Paper Tokens & Skill Tokens & Skill/Full & Reduction \\
\midrule
AI Scientist-v2 & 45212 & 1079 & 0.0239 & 97.61\% \\
Reflexion & 16414 & 703 & 0.0428 & 95.72\% \\
AIDE & 13312 & 1285 & 0.0965 & 90.35\% \\
Toolformer & 20365 & 1255 & 0.0616 & 93.84\% \\
\bottomrule
\end{tabular}
\end{table*}

\begin{table*}[t]
\centering
\caption{Automatic extracted-text note scaffolds are evaluated separately from curated-note cases.}
\label{tab:auto-note}
\begin{tabular}{llrrrrr}
\toprule
Paper & Input & Rubric & Coverage & Transfer & Support & Words \\
\midrule
Toolformer & Curated note & 20/20 & 8.9/10 & 10/10 & 1.000 & 943 \\
Toolformer & Auto note scaffold & 20/20 & 9.3/10 & 10/10 & 1.000 & 1179 \\
AIDE & Curated note & 20/20 & 9.1/10 & 10/10 & 1.000 & 927 \\
AIDE & Auto note scaffold & 20/20 & 8.467/10 & 9.5/10 & 1.000 & 998 \\
\bottomrule
\end{tabular}
\end{table*}


\section{Results}

The completed results below are deterministic/offline quality and grounding
evidence for the current PaperToSkill artifacts, plus first-pass real-reuse
execution evidence. AIDE-T1 and AIDE-T2 are scored for one GPT-family run on a
locked Kaggle Spaceship Titanic validation split using the same saved model
outputs and an extended 300-second local scorer budget. AIDE-T1 is solved by
both Summary and PaperToSkill with nearly tied validation scores (0.816 and
0.817), while AIDE-T2 is a PaperToSkill-only success (0.826 versus a Summary
timeout). SWE-T1 is scored for one GPT-family run on a locked
SWE-Bench Lite SQLFluff instance: both Summary and PaperToSkill produce patches
that fail to apply and score 0.000. SWE-T2 is scored for one GPT-family run on
a locked SWE-Bench Verified-style Astropy instance: the Summary patch fails to
apply and scores 0.000, while the PaperToSkill patch applies, passes the hidden
target tests, and scores 1.000. In the real-reuse table, REF-T1 and REF-T2 are
also scored for one GPT-family run: both Summary and PaperToSkill reach 1.000
on the locked Reflexion QA and retry fixtures. SNAP-T1 and SNAP-T2 are scored
for one GPT-family run over prepared official miniature fixtures. PaperToSkill
scores above Summary on these two SnapATAC2 dry-run tasks (0.500 vs. 0.000 and
0.400 vs. 0.200), but all SNAP scores remain below the success threshold,
mainly because the responses did not produce complete runtime, memory, and
quality artifacts. Across the eight filled tasks, the runner, scorer, and
table-update paths are live. The outcome is mixed: AIDE-T2 and SWE-T2 are
positive PaperToSkill rows, AIDE-T1 and REF show little or no advantage because
both conditions solve the fixtures, and SWE-T1/SNAP primarily expose
patch-application and artifact-completion failure boundaries. These results do
not establish
aggregate advantage over Summary. Table~\ref{tab:real-reuse-failure-analysis}
summarizes the row-level boundary modes and the method contracts they suggest
for follow-up work. The SWE-T1 diagnostics separate two contract issues. The
shared-source-context follow-up
(Table~\ref{tab:swe-t1-source-context-followup}) fixes patch application for
both conditions but still fails the hidden target test. The issue-aligned
follow-up (Table~\ref{tab:swe-t1-issue-aligned-followup}) then uses a
pre-registered revised scorer and both Summary and PaperToSkill pass, so it
shows that the revised task contract can close rather than showing a
PaperToSkill advantage. A paired SnapATAC2 executable-artifact diagnostic
(Table~\ref{tab:snapatac2-executable-followup}) then runs a pre-registered
controlled scaffold over the same miniature fixtures. It reaches the SNAP
scorer's artifact/runtime/memory contract for both Summary and PaperToSkill
conditions, confirming that the prior SNAP failure mode is an execution
contract boundary rather than a fixture or scorer impossibility; it is not a
main-row replacement and does not show a PaperToSkill advantage. A subsequent
model-generated executable-candidate diagnostic for SNAP-T1
(Table~\ref{tab:snapatac2-executable-candidate-followup}) also reaches 1.000
for both Summary and PaperToSkill under the pre-registered executable-candidate
contract. This closes the candidate-script artifact contract for both
conditions; it does not replace the main SNAP rows or show a PaperToSkill
advantage. The auxiliary Full Excerpt sanity subset does not reverse this
interpretation: AIDE-T1 and SWE-T1 score 0.000 under Full Excerpt, while
SNAP-T1 scores 0.250, below both PaperToSkill and the success threshold.
All four generated skills score 20/20 on deterministic skill rubrics. On context
coverage, the generated skills score 7.867/9 for AI Scientist-v2, 8.267/9 for
Reflexion, 9.1/10 for AIDE, and 8.9/10 for Toolformer. The generic-summary and
abstract-only baselines score substantially lower across all four cases
(Table~\ref{tab:main-results}).

The generated skills remain compact: 782 words for AI Scientist-v2, 479 words
for Reflexion, 927 words for AIDE, and 943 words for Toolformer, all under the
1200-word budget. Source validation finds support rates of 0.938, 1.0, 1.0,
and 1.0 with zero invalid line ranges. The one weak AI Scientist-v2 span is
recorded as a boundary case rather than treated as fully verified.

The tokenizer-aware context proxy shows the same compactness pattern relative
to full paper contexts. Under \texttt{o200k\_base}, generated skills use between
2.39\% and 9.65\% of the input tokens required by full extracted paper text,
with the full tokenizer table retained in the package artifacts. A separate
character-proxy report preserves the earlier
$\lceil\mathrm{characters}/4\rceil$ estimate for reproducibility. We interpret
both as coverage-preserving compression relative to full-paper context, not as
provider billing evidence or a claim that skills are the shortest possible
context.

Transfer-note ablations show a consistent offline readiness pattern. Full
skills score 10/10 on the readiness metric across all four curated papers.
Removing \texttt{Transfer Notes} lowers readiness to 7.6/10 in all four cases,
while generic summaries score between 1.2/10 and 2.25/10. This supports the
narrower claim that transfer notes encode artifact-readiness signals. It does
not yet prove improved live cross-harness success.

The live-transfer saved-response evaluation covers all four paper packets.
AI Scientist-v2, Reflexion, AIDE, and Toolformer each have six saved responses
across Codex-style and Claude-style harness prompts and three context variants.
The aggregate evaluator reports 24 total rows, 24 scored rows, 0 pending rows,
and average normalized score 1.0. AI Scientist-v2, Reflexion, and AIDE rows
score 11/11 each; Toolformer rows score 9/9 each. The AIDE run includes one
provider fallback row where \texttt{claude-opus-4-8} closed the connection and
\texttt{claude-opus-4-7} succeeded. These are deterministic output-contract
scores over saved response files, not human-rated semantic fidelity or
downstream execution-outcome rates.

The automatic note scaffold is evaluated separately on Toolformer and AIDE. The
Toolformer auto-note-derived skill scores 20/20 on the deterministic rubric,
9.3/10 on context coverage, 10/10 on offline transfer readiness, and 1.0
source-span support with zero invalid ranges. The AIDE auto-note-derived skill
scores 20/20, 8.467/10 on context coverage, 9.5/10 on offline transfer
readiness, and 1.0 source-span support with zero invalid ranges. The full
auto-note comparison table is retained in the package artifacts. These results
support the narrower claim that extracted text can seed auditable note
scaffolds. The local pipeline also
includes a direct-PDF smoke path through \texttt{pdftotext -layout} that records
the generated text source in the manifest, but this does not establish robust
arbitrary-PDF automation.

We additionally compare PaperToSkill against Paper2Agent at the artifact and
workflow level using the Paper2Agent paper text and current PaperToSkill
artifacts. The bounded comparison covers required inputs, generated artifact
type, setup burden, validation checks, failure handling, source traceability,
and runtime dependency. All seven criteria are source-backed and ready in the
local comparison table. The result clarifies positioning: Paper2Agent produces
MCP servers for papers with associated codebases, while PaperToSkill produces
portable natural-language skills with explicit source and failure boundaries.
This comparison does not run Paper2Agent, deploy an MCP server, or establish an
end-to-end baseline performance result.

\section{Usage Examples and Model Ablations}

The repository includes usage examples for three practical workflows:
loading a generated skill in a Codex-style harness, generating an automatic
note scaffold from extracted text or a smoke-tested PDF source, and building
model-ablation prompts. The
model-ablation protocol produces a prompt grid for Claude Opus 4.8, a
GPT-family slot requested as GPT 5.5, and a DeepSeek slot. In the current
protocol, the saved-response evaluator reports 6 total rows, 6 scored rows, 0
pending rows, and average normalized score 1.0. The GPT-family protocol refresh
completed both rows with \texttt{gpt-5.5}. The DeepSeek run completed both rows
with \texttt{deepseek-v4-flash}. The latest Claude protocol refresh used the
correct Anthropic Messages path but was blocked by provider HTTP 502; the
scored Claude rows come from previously saved response files. A local
output-token proxy over all six saved model-ablation responses counts 9,594
\texttt{o200k\_base} output tokens. This 6/6 saved-response score is not live
downstream task success, provider billing, live-invoice, success-per-dollar
evidence, or a broad model-quality result.

Separately, the real-reuse LLM ablation attaches model variation to the
paper-task protocol rather than to the older saved-response usage-plan
protocol (Table~\ref{tab:real-reuse-llm-ablation}). The pre-registered slice
uses AIDE-T2, SWE-T2, and REF-T2 under the same Summary/PaperToSkill pairing.
GPT-family and DeepSeek-family rows are collected, while Claude-family rows
remain pending after provider HTTP 502 responses. The collected rows are mixed:
GPT-family favors Summary on AIDE-T2 and ties on REF-T2, while both conditions
fail SWE-T2; DeepSeek ties on the collected slices. This table is auxiliary
model-slice evidence and does not replace the main real-reuse rows.

The bounded AI-Scientist-v2 integration path is complete for the local marker
smoke and one full live run. The smoke report records a complete marker-contract
response, and the full-run handoff records one completion directory. The
AI-Scientist-v2-generated synthetic benchmark ties skill and full-excerpt
task-success rates at 0.80 while the skill uses fewer tokens (86.2 versus
113.2); abstract, generic-summary, and no-context baselines score 0.20, 0.00,
and 0.00. A retrieval-depth sensitivity branch reaches 1.00 only when K=all.
This is bounded integration and synthetic sensitivity evidence, not human
fidelity, real-data validation, or broad live research-task success.

\section{Discussion}

The main result is not that PaperToSkill ``understands'' papers in the broadest
sense. The result is that a structured paper-note-to-skill pipeline can preserve
operational knowledge that summaries discard. The skills include workflow
steps, validation checks, failure cases, and source anchors, which are exactly
the parts needed when a user wants an agent to apply a paper rather than merely
explain it.

The benchmark also reveals useful failure modes. During the AIDE case, the
initial extractor capped workflow, validation, and failure candidates too
aggressively, which dropped data-preview and LLM-cost content. Earlier phases
also exposed title parsing, source-audit mapping, and source-span line-counting
bugs. These records are the type of failure branches PaperToSkill should
preserve: not just final successes, but the ways extraction, evaluation, or
external dependencies can lose important operational content.

The real-reuse rows add a harsher version of the same lesson. The AIDE rows
show that task slice and local scorer budget can change the apparent boundary:
with a longer scorer budget AIDE-T1 is solved by both conditions and AIDE-T2
becomes a PaperToSkill-only success, while the Summary AIDE-T2 timeout still
argues for explicit runtime and fallback contracts. SWE-T1 shows that
software-engineering reuse can fail before tests run if a patch cannot be
applied cleanly. The SnapATAC2 rows show that partial pipelines are not enough
when the scorer requires complete runtime, memory, and quality artifacts. These
results point to the next method pressure points: budget contracts,
patch/application constraints, required artifact manifests, and recovery
instructions should be explicit when a source paper's method is converted into
a reusable skill.

\section{Limitations}

The current work has several important limits. First, the main benchmark
converts curated notes rather than arbitrary PDFs; automatic note scaffolds have
only been retained on Toolformer and AIDE extracted text. Second, the metrics
are deterministic and partly lexical, so they can over-credit exact matches or
under-credit valid paraphrases. Third, live-transfer saved-response coverage is
complete for the current prompt-packet protocol, but the scorer is
deterministic and contract-based rather than a human semantic-fidelity or live
execution-outcome evaluation. Fourth, the bounded AI-Scientist-v2 smoke and
full live run is complete, but its positive result is synthetic and should not
be read as broad live task success. Fifth,
human-fidelity packets, a reviewer handoff guide, a stricter blank annotation
template, and a summary script are prepared, but no independent annotations have
been completed. Sixth,
compactness is measured by word count, deterministic character proxy, local
tokenizer-aware input-token proxy, and local output-token proxy for saved model
responses, not by provider-specific prices, live invoices, or live success per
dollar. Seventh, the real-reuse rows are single-run GPT-family measurements,
and original-paper reference scores are reported references unless the same
data, input/output contract, metric, budget, and runtime setting are reproduced
locally. Several rows fail because generated scripts or patches do not finish
within the locked budget or do not produce complete scorer artifacts, so they
should be read as boundary evidence rather than a broad downstream
effectiveness claim.

\section{Conclusion}

PaperToSkill turns paper-derived procedural knowledge into portable,
human-editable skills. In a four-paper benchmark, generated skills are compact,
source-grounded, structurally valid, and more operationally complete than short
summary baselines under deterministic evaluation. A first eight-row real-reuse
stress test is now complete, but it is mixed and failure-heavy rather than a
proof of aggregate downstream effectiveness. The next stage is to add human
fidelity review, stabilize the locked real-reuse rows with pre-registered
paired follow-ups where needed, extend the bounded Paper2Agent artifact
comparison into a real executable MCP baseline if resources permit, and use
the real-reuse failure modes to refine budget, artifact, and failure-recovery
contracts.

\bibliography{papertoskill_refs}

\end{document}
